\section{Canonistic Formation, Scholastic Synthesis, and Early Modern Expansion}

\subsection{Why systematic canon law emerged}

Before Gratian, the Latin Church possessed Scripture, patristic authorities, conciliar canons, papal decretals, penitential norms, liturgical sources, and regional collections. Their dates, addressees, purposes, and formulations differed. Eleventh- and twelfth-century reform disputes over simony, clerical life, episcopal appointment, papal authority, and lay control increased the need to identify and reconcile governing sources. At the same time, renewed study of Justinianic Roman law and the growth of schools supplied juristic methods.

This conjunction matters. Systematic canon law did not arise only from an abstract interest in classification. Expanding ecclesiastical jurisdiction, reform, litigation, legislation, and teaching created practical questions about which authority and text controlled. The recovery of Roman law contributed vocabulary and procedure; ecclesial sources supplied a distinct constitution and end. Anders Winroth's account connects reform, papal government, the schools, and increasingly systematic collections while preserving the contested chronology of Gratian's work (\sourceurl{https://www.cambridge.org/core/services/aop-cambridge-core/content/view/FA2BFB47AB644DA069661A47E69B08D9/9781139056021c6_p89-98_CBO.pdf/the-legal-underpinnings.pdf}{historical control}).

\subsection{Gratian: discordant authorities and juristic concord}

The \emph{Decretum} begins by saying that humankind is ruled by natural law and customs. Gratian identifies natural law with what is contained in the Law and Gospel, including the command to treat another as one wishes to be treated; his broad category of custom includes written and unwritten human law (D.1 dictum before c.1; \sourceurl{https://geschichte.digitale-sammlungen.de/decretum-gratiani/kapitel/dc_chapter_0_0004}{Friedberg edition text}). D.1 c.1, drawn from Isidore's \emph{Etymologies} V.2, then divides laws into divine, grounded in nature, and human, grounded in custom (\sourceurl{https://geschichte.digitale-sammlungen.de/decretum-gratiani/kapitel/dc_chapter_0_5}{D.1 c.1}).

The opening contains real tension. Gratian's initial identification is scriptural; D.1 cc.6--9 also transmit Isidore's Roman-derived natural, civil, and gentium categories (\sourceurl{https://geschichte.digitale-sammlungen.de/decretum-gratiani/kapitel/dc_chapter_0_10}{D.1 c.6}; \sourceurl{https://geschichte.digitale-sammlungen.de/decretum-gratiani/kapitel/dc_chapter_0_11}{c.7}; \sourceurl{https://geschichte.digitale-sammlungen.de/decretum-gratiani/kapitel/dc_chapter_0_12}{c.8}; \sourceurl{https://geschichte.digitale-sammlungen.de/decretum-gratiani/kapitel/dc_chapter_0_13}{c.9}). Gratian's distinctions, hypothetical cases, dicta, and reconciliation techniques offered later decretists a method for determining whether authorities addressed different times, places, meanings, or exceptions.

D.4 c.2 receives criteria for human enactment: it should be honorable, just, possible, fitting to nature and local custom, adapted to time and place, necessary, useful, clear, and directed to common utility (\sourceurl{https://geschichte.digitale-sammlungen.de/decretum-gratiani/kapitel/dc_chapter_0_33}{text}). Juridical form and substantive justice are both part of the inquiry.

\subsection{Decretists, decretals, courts, and the \emph{ius commune}}

Rufinus's \emph{Summa Decretorum} distinguishes a natural capacity implanted in rational humans from the Roman animal-wide use of \emph{ius naturale}. He relates institutions understood to arise after the Fall, the \emph{ius gentium}, the Decalogue and Gospel, and positive canons governing ecclesial institutions (ed. Heinrich Singer, 1902, pp.~4--10; \sourceurl{https://archive.org/details/diesummadecretor00rufi}{scan}). His work exemplifies decretist efforts to interpret, rather than merely repeat, Gratian's authorities.

Papal decretals issued in cases and as legislation continued to accumulate. Gregory IX's \emph{Liber Extra}, compiled by Raymond of Peñafort and promulgated in 1234, reorganized this material into five books for teaching and practice (\sourceurl{https://web.colby.edu/canonlaw/extra/}{edition and source guide}). Decretalists interpreted the new collection; later official and private collections joined Gratian and the \emph{Liber Extra} in what became the \emph{Corpus iuris canonici}.

The medieval \emph{ius commune} was not one statute book. The learned \emph{ius commune} comprised Roman and canon law and their juristic literature; courts applied it in interaction with local, feudal, municipal, corporate, royal, and customary \emph{iura propria}. Ecclesiastical and civil courts could apply overlapping bodies of doctrine to marriage, oath, property, status, contract, benefice, crime, and procedure. A litigant's forum, status, cause of action, proof, and remedy affected the outcome. Modern language of legal pluralism helps describe that multiplicity, provided it does not project fully separate modern legal systems backward onto every source and institution.

Romano-canonical procedure developed structured pleading, proof, representation, appeal, and record. These forms did more than administer substantive norms: they determined what could be established and remedied in the external forum. Claims of right or natural justice could receive enforceable effect when recognized through actions, exceptions, presumptions, and competent judgment.

\subsection{Penitential and contentious forums}

Medieval Christians also encountered norms through penance and confession. Canonists and pastoral writers distinguished forms of judgment concerned with conscience, sacramental discipline, public status, and contentious litigation. The vocabulary and institutional boundary developed over time. Joseph Goering traces a two-fora account in the decades after Gratian while warning that later language of the \emph{forum internum} should not be projected indiscriminately into earlier texts (\sourceurl{https://www.cambridge.org/core/journals/traditio/article/abs/internal-forum-and-the-literature-of-penance-and-confession/D2C46A52A2AD91AE721EA5162AEA203D}{study}). Wolfgang P. Müller's contrary historiographical challenge shows that the emergence and later reconstruction of the category remain disputed (\sourceurl{https://doi.org/10.1017/S0738248015000486}{study}).

The historical distinction matters for personal obligation. A judgment about sin, an ecclesiastical penalty, a sacramental remedy, and a civil cause of action could concern related conduct without producing identical effects. Current canon law's internal and external forums have positive definitions and cannot simply be read back into medieval practice.

\subsection{Aquinas's architecture}

Aquinas relates law to beatitude, virtue, grace, prudence, community, and salvation in \emph{Summa theologiae} I--II qq.~90--108:

\begin{fourstable}{Form}{Source and character}{Mode of knowledge}{Function}
Eternal law & Divine wisdom governing the created order. & Known perfectly by God; creatures participate according to their mode. & Ultimate source and measure of just order.\\
Natural law & Rational participation in eternal law. & Practical reason grasps first principles and reasons toward conclusions. & Directs action toward human goods and supplies human law's moral basis.\\
Human law & Particular conclusions and determinations for a political community. & Known through promulgated enactment or legally operative custom and interpreted by reason and authority. & Coordinates action, restrains grave wrongdoing, secures justice, and supports virtue.\\
Divine law & The revealed Old and New Law; the New Law is principally grace. & Received through revelation and faith. & Directs persons to their supernatural end and reaches acts beyond human law's competence.\\
\end{fourstable}

The principal loci are \sourceurl{https://www.newadvent.org/summa/2090.htm}{q.~90}, \sourceurl{https://www.newadvent.org/summa/2091.htm}{q.~91}, \sourceurl{https://www.newadvent.org/summa/2093.htm}{q.~93}, \sourceurl{https://www.newadvent.org/summa/2094.htm}{q.~94}, and \sourceurl{https://www.newadvent.org/summa/2095.htm}{q.~95}; Old and New Law receive extended treatment in qq.~98--108.

Aquinas also treats the stability and change of human law, custom, dispensation, and equity. Change can serve the common good while weakening law's habituating stability (q.~97, a.~2). Custom can acquire or interpret legal force through the community's juridically relevant action (q.~97, a.~3). Dispensation belongs to authority entrusted with the community and applies a common rule where its reason fails in the case (q.~97, a.~4). Equity addresses exceptional application of general wording according to the rule's rational purpose (II--II q.~120). Current positive law defines which authorities and procedures can give those concerns external legal effect.

\subsection{Salamanca, conquest, and the law of peoples}

European expansion forced scholastic and juridical categories into disputes about the status, property, political authority, religious freedom, and resistance of non-Christian peoples. In \emph{De Indis}, Francisco de Vitoria tests asserted Spanish titles to dominion against natural law and the \emph{ius gentium} and recognizes the political and proprietary \emph{dominium} of Indigenous peoples (``On the American Indians,'' q.~1, conclusion, pp.~239--251 in the Pagden--Lawrance edition). His arguments also retained contested titles and assumptions connected to communication, mission, and war; they should not be presented as a completed anticolonial settlement.

The Valladolid controversy of 1550--1551 placed Juan Ginés de Sepúlveda's arguments about conquest and hierarchy against Bartolomé de las Casas's defense of Indigenous rationality, freedom, and evangelization without coercive war. The sources disclose a live conflict over jurisdiction, cultural difference, natural rights, religious mission, and imperial power rather than a linear triumph of natural law (\sourceurl{https://academic.oup.com/book/46077}{documentary and historiographical guide}). The episode demonstrates both the critical capacity of higher-law reasoning and its vulnerability to disputed anthropology and imperial interests.

\subsection{Trent, Suárez, and the early modern bridge}

The Council of Trent supplies concrete examples of ecclesiastical competence and positive determination. Session VII, canon 8 on baptism rejects freedom from all Church precepts. Session XXIV, canon 4 on matrimony affirms ecclesiastical competence to establish diriment impediments; \emph{Tametsi} then enacts a form requirement where promulgated (\sourceurl{https://www.papalencyclicals.net/councils/trent/seventh-session.htm}{Session VII}; \sourceurl{https://www.papalencyclicals.net/councils/trent/twenty-fourth-session.htm}{Session XXIV witnesses}). The doctrinal claim about competence and the changeable disciplinary rule remain distinct.

Francisco Suárez's \emph{De legibus ac Deo legislatore} (1612) examines law and lawgiver, obligation and promulgation, eternal and natural law, human civil and canonical law, \emph{ius gentium}, custom, privilege, and revealed law (especially I.5--12; II; III--IV; VII--X). He locates the legislator's will within rational direction of a community toward the common good. His account is an early modern scholastic synthesis with its own terminology; it should not be reduced either to Thomistic repetition or to later command theory.

Reformation, confessional division, overseas empire, stronger territorial government, and disputes over sovereignty changed the institutions within which these categories operated. Natural law and \emph{ius gentium} became languages used across confessional boundaries, while civil and ecclesiastical jurisdictions were increasingly reorganized by states. The next section follows that transformation into toleration, codification, legal positivism, rights, and Vatican II.
